\documentclass[11pt,a4paper]{article}
\usepackage[a4paper,margin=2.2cm]{geometry}
\usepackage{fontspec}
\setmainfont{Times New Roman}
\setmonofont{Consolas}[Scale=0.86]
\usepackage{amsmath,amssymb,mathtools}
\usepackage{unicode-math}
\setmathfont{Latin Modern Math}
\usepackage{booktabs}
\usepackage{array}
\usepackage{enumitem}
\usepackage{titlesec}
\usepackage{xcolor}
\usepackage{mdframed}
\usepackage[hidelinks]{hyperref}

\definecolor{acc}{HTML}{7A2E1E}
\definecolor{dim}{HTML}{555555}
\definecolor{boxbg}{HTML}{F4F1EC}

\titleformat{\section}{\large\bfseries\color{acc}}{\thesection.}{0.5em}{}
\titlespacing{\section}{0pt}{1.1em}{0.4em}
\titleformat{\subsection}{\bfseries}{\thesubsection.}{0.5em}{}
\titlespacing{\subsection}{0pt}{0.8em}{0.25em}

\setlist[itemize]{leftmargin=1.2em,itemsep=0.15em,topsep=0.3em}
\setlist[enumerate]{leftmargin=1.4em,itemsep=0.15em,topsep=0.3em}
\renewcommand{\arraystretch}{1.15}
\setlength{\parskip}{0.45em}
\setlength{\parindent}{0pt}

\newcommand{\qq}[1]{\texttt{\small #1}}
\newcommand{\note}[1]{{\color{dim}\small #1}}

% Врезка для вывода или ключевого следствия.
\newmdenv[
  backgroundcolor=boxbg, linewidth=0pt, skipabove=0.7em, skipbelow=0.7em,
  innerleftmargin=0.9em, innerrightmargin=0.9em,
  innertopmargin=0.6em, innerbottommargin=0.6em
]{keybox}

\DeclareMathOperator*{\argmax}{arg\,max}
\DeclareMathOperator*{\argmin}{arg\,min}
\DeclareMathOperator{\KL}{KL}
\newcommand{\R}{\mathbb{R}}
\newcommand{\E}{\mathbb{E}}
\newcommand{\Var}{\operatorname{Var}}

\begin{document}

{\Large\bfseries Дообучение: математика и методы}

\vspace{-0.5em}
\note{vlm-recipes~· подзадача 2}

\section{Что именно оптимизируется}

Языковая модель задаёт распределение над последовательностями токенов
через авторегрессионную факторизацию:
\[
p_\theta(x_1,\dots,x_T) \;=\; \prod_{t=1}^{T} p_\theta(x_t \mid x_{<t}).
\]
Обучение — минимизация отрицательного логарифма правдоподобия, то есть
кросс-энтропии между предсказанным распределением и наблюдаемым токеном:
\[
\mathcal{L}(\theta) \;=\; -\sum_{t=1}^{T} \log p_\theta(x_t \mid x_{<t}).
\]

Здесь и возникает первая практическая развилка. Сумма идёт по всем
позициям, включая промпт пользователя, системную инструкцию и результаты
вызовов инструментов. Модель, обученная так, учится в том числе
\emph{порождать} вопросы пользователя и \emph{выдумывать} вывод
инструментов.

Маскирование ограничивает сумму множеством $A$ позиций, принадлежащих
репликам ассистента:
\[
\mathcal{L}_{\text{SFT}}(\theta) \;=\; -\frac{1}{|A|}\sum_{t \in A}
\log p_\theta(x_t \mid x_{<t}).
\]

\begin{keybox}
Существенно, что $x_{<t}$ по-прежнему включает \emph{весь} контекст.
Маскирование убирает позиции из функции потерь, но не из обусловливания:
модель видит промпт, просто не получает градиента за его предсказание.
Технически это реализуется значением $-100$ в метках, которое
\qq{cross\_entropy} пропускает.
\end{keybox}

Отсюда же диагностический признак: если доля обучаемых токенов
$|A|/T$ мала (единицы процентов), почти весь вычислительный бюджет
батча уходит на прямой проход по тексту, за который нет градиента.

\section{Полное дообучение и его цена}

Обновляются все $P$ параметров. Помимо самих весов в памяти живут
градиенты и состояния оптимизатора. Для AdamW со смешанной точностью:

\begin{center}
\begin{tabular}{@{}lr@{}}
\toprule
веса bf16 & $2P$ \\
мастер-копия fp32 & $4P$ \\
градиенты & $2P$ \\
моменты AdamW $m,v$ (fp32) & $8P$ \\
\midrule
итого & $\approx 16P$ байт \\
\bottomrule
\end{tabular}
\end{center}

Для $P = 9\cdot10^9$ это около 144 ГБ — без учёта активаций. Отсюда
и растут все остальные методы: они атакуют разные слагаемые этой суммы.

\section{LoRA: низкоранговая адаптация}

\subsection{Конструкция}

Идея опирается на наблюдение Aghajanyan и др. о низкой внутренней
размерности задач дообучения: приращение весов $\Delta W$, нужное для
адаптации к узкой задаче, хорошо приближается матрицей малого ранга.

Для слоя с весами $W_0 \in \R^{d \times k}$ приращение параметризуется
произведением
\[
\Delta W = BA, \qquad B \in \R^{d \times r},\; A \in \R^{r \times k},
\qquad r \ll \min(d,k),
\]
а прямой проход становится
\[
h = W_0 x + \gamma\, BAx, \qquad \gamma = \frac{\alpha}{r}.
\]

Число обучаемых параметров падает с $dk$ до $r(d+k)$. При $d=k=4096$
и $r=32$ это $2{,}6\cdot10^5$ вместо $1{,}7\cdot10^7$ — в шестьдесят
четыре раза меньше.

\subsection{Инициализация}

$A$ инициализируется случайно, $B$ — нулём. Тогда в начале обучения
$\Delta W = 0$, и модель стартует ровно из претрейна: адаптер ничего
не портит до первого шага.

Возникает вопрос, почему не наоборот. Выпишем градиенты, обозначив
$g = \partial L/\partial h$:
\[
\frac{\partial L}{\partial B} = \gamma\, g\,(Ax)^\top,
\qquad
\frac{\partial L}{\partial A} = \gamma\, B^\top g\, x^\top .
\]

При $A$ случайном и $B=0$ градиент по $A$ обнуляется множителем
$B^\top$, но градиент по $B$ содержит $Ax \neq 0$ — значит $B$
сдвигается на первом же шаге, после чего оживает и $A$. При $A=0$
и $B$ случайном картина зеркальна: сдвигается $A$, следом оживает $B$.

\begin{keybox}
Обе схемы обучаются. Сломан только третий случай — обе матрицы
нулевые: тогда оба градиента нулевые навсегда. Ноль нужен ровно
в одной из матриц, чтобы дать $\Delta W = 0$ на старте, и случайность
в другой, чтобы разорвать симметрию.
\end{keybox}

Выбор между двумя рабочими схемами определяется не запуском обучения,
а его динамикой. Hayou и др. (2024) показывают, что схемы не
эквивалентны: при случайном $A$ и нулевом $B$ допустима бо́льшая
скорость обучения без потери устойчивости. Интуиция в том, что на
первом шаге $A$ работает как фиксированная случайная проекция входа,
а $B$ учится уже в спроецированном пространстве; размерности матриц
различны, поэтому различны и масштабы обновлений.

\subsection{Почему память экономится в разы}

Замороженные $W_0$ не требуют ни градиентов, ни моментов. Состояния
оптимизатора нужны только для $BA$, то есть для доли процента
параметров. Веса остаются в памяти (их нужно читать на прямом проходе),
но три остальных слагаемых из таблицы выше исчезают почти целиком.

\section{Масштаб: почему \texorpdfstring{$\alpha/r$}{alpha/r} плох при больших рангах}

Это место стоит разобрать, потому что оно объясняет распространённое
наблюдение «ранг выше 32 ничего не даёт».

Пусть после нескольких шагов элементы $A$ и $B$ имеют порядок
$\sigma_A$ и $\sigma_B$ и приблизительно независимы. Элемент
произведения есть сумма $r$ слагаемых:
\[
(BA)_{ij} = \sum_{s=1}^{r} B_{is} A_{sj},
\qquad
\Var\big[(BA)_{ij}\big] \approx r\,\sigma_B^2 \sigma_A^2 .
\]
Среднеквадратичная величина элемента растёт как $\sqrt{r}$. С множителем
$\gamma = \alpha/r$ вклад адаптера в выход имеет масштаб
\[
\gamma \cdot \sqrt{r}\,\sigma_B\sigma_A
= \frac{\alpha\,\sigma_B\sigma_A}{\sqrt{r}} \;\xrightarrow[r\to\infty]{}\; 0 .
\]

\begin{keybox}
При $\gamma=\alpha/r$ вклад адаптера \emph{затухает} как $1/\sqrt{r}$.
Увеличивая ранг, вы одновременно глушите адаптер, и два эффекта гасят
друг друга. Правильный масштаб —
\[
\gamma = \frac{\alpha}{\sqrt{r}},
\]
при нём величина вклада от $r$ не зависит. Это и есть \textbf{rsLoRA}
(rank-stabilized LoRA, Kalajdzievski, 2023); в PEFT включается флагом
\qq{use\_rslora=True}.
\end{keybox}

Практический вывод: при $\gamma=\alpha/r$ разумный диапазон рангов
кончается около 32. С rsLoRA ранги 64 и 128 начинают приносить реальный
прирост. Стоит это ничего, поэтому включать следует всегда.

\section{DoRA: разделение величины и направления}

Эмпирическое наблюдение Liu и др. (2024): при полном дообучении норма
столбца весов и его направление меняются \emph{независимо} — бывают
большие повороты при почти неизменной норме и наоборот. LoRA такой
свободы не даёт: одно приращение $BA$ управляет обоими сразу.

DoRA разлагает веса явно:
\[
W = m \odot \frac{V}{\lVert V \rVert_c},
\]
где $\lVert\cdot\rVert_c$ — норма по столбцам, $m \in \R^{1\times k}$ —
обучаемый вектор величин, $V$ — направление. Дообучение применяет
низкоранговое приращение только к направлению:
\[
W' = m \odot \frac{W_0 + BA}{\lVert W_0 + BA \rVert_c}.
\]

Обучаются $m$, $A$, $B$. Нормировка развязывает две степени свободы,
и на низких рангах ($r \le 16$) DoRA устойчиво обходит LoRA. Цена —
вычисление нормы на каждом шаге, около 20\% скорости.

\section{Квантование базы: NF4 и QLoRA}

\subsection{Почему не обычный int4}

Равномерная сетка из 16 уровней плоха для весов нейросети, потому что
они распределены приблизительно нормально: большинство значений
сосредоточено у нуля, а равномерные уровни тратят разрешение на хвосты.

\textbf{NF4} (NormalFloat4) строит уровни как квантили стандартного
нормального распределения. Уровни $q_i$ выбираются так, чтобы каждый
интервал между ними нёс равную вероятностную массу:
\[
q_i \;\propto\; Q_{\mathcal{N}}\!\left(\frac{i + 0.5}{2^b}\right),
\qquad i = 0,\dots,2^b-1,
\]
где $Q_{\mathcal{N}}$ — квантильная функция. Для входа, распределённого
нормально, такая сетка информационно-теоретически оптимальна.

\subsection{Блоки и двойное квантование}

Веса нормируются поблочно: блок из 64 значений делится на свой масштаб
$c = \max|w|$, и уже отнормированное квантуется в NF4. Масштабы
хранятся отдельно и сами занимают память: при блоке 64 и fp32-масштабе
это $32/64 = 0{,}5$ бита на параметр.

\textbf{Двойное квантование} сжимает и масштабы: они квантуются в 8 бит
блоками по 256, что снижает накладные расходы до $\approx 0{,}127$ бита
на параметр. Итого QLoRA хранит вес примерно в $4{,}127$ бита.

\subsection{Что происходит на прямом проходе}

\[
h = \text{dequant}_{\text{bf16}}(W^{\text{NF4}})\,x + \gamma BAx .
\]

Умножение выполняется в bf16: веса разжимаются поблочно на лету.

\begin{keybox}
Отсюда важное следствие: \textbf{QLoRA экономит память, но не ускоряет}.
Разжатие добавляет работу, поэтому шаг обучения медленнее, чем у LoRA
в bf16. Брать QLoRA следует тогда и только тогда, когда bf16-база
не помещается.
\end{keybox}

Обучаемые $A$, $B$ остаются в bf16 — квантованию подвергается только
замороженная база, иначе градиенты потеряли бы разрешение.

\section{Ортогональные приёмы}

\subsection{NEFTune}

К эмбеддингам входа при обучении добавляется равномерный шум:
\[
x'_{\text{emb}} = x_{\text{emb}} + \varepsilon,
\qquad
\varepsilon \sim \mathcal{U}\!\left(-\frac{\alpha}{\sqrt{Ld}},\,
\frac{\alpha}{\sqrt{Ld}}\right),
\]
где $L$ — длина последовательности, $d$ — размерность эмбеддинга.
Нормировка на $\sqrt{Ld}$ делает величину возмущения независимой
от длины примера.

Эффект — регуляризация: модель перестаёт цепляться за точную форму
обучающих ответов. На малых выборках прирост заметный, стоимость
нулевая. Обычное значение $\alpha = 5$.

\subsection{Слитые ядра (Liger)}

Наивное вычисление лосса материализует логиты размера
$B \times L \times |V|$. При словаре $|V| \sim 150{,}000$ это доминирующее
слагаемое в памяти активаций. Слитые ядра считают проекцию в словарь
и кросс-энтропию чанками, не храня полную матрицу логитов. Экономия тем
больше, чем крупнее словарь; на качество не влияет никак.

\subsection{GaLore: низкий ранг в градиенте, а не в весах}

Альтернатива LoRA, сохраняющая полноранговое обучение весов.
Градиент $G_t \in \R^{m\times n}$ проецируется в низкоранговое
подпространство, и состояния оптимизатора живут уже там:
\[
\tilde{G}_t = P^\top G_t, \qquad P \in \R^{m \times r},
\]
где $P$ — первые $r$ левых сингулярных векторов $G_t$, пересчитываемые
раз в несколько сотен шагов. Обновление возвращается в исходное
пространство: $\Delta W = -\eta\, P\,\rho(\tilde{G}_t)$.

Смысл в том, что ограничение накладывается на оптимизатор, а не на
выразительность модели: веса меняются полноранговым образом. Память
моментов падает в $m/r$ раз.

\section{Что выбрать}

\begin{center}
\begin{tabular}{@{}lp{5.4cm}p{5.4cm}@{}}
\toprule
\textbf{Метод} & \textbf{Когда} & \textbf{Цена} \\
\midrule
full & максимум качества, ресурсы есть & $16P$ байт памяти \\
lora & базовый выбор & $r$ выше 32 бесполезен \\
rslora & \textbf{всегда вместо lora} & нет \\
dora & низкие ранги, тонкая настройка стиля & $-20\%$ скорости \\
qlora & база не помещается в bf16 & медленнее lora \\
qdora & то же плюс низкий ранг & медленнее всех \\
galore & нужно полноранговое обучение при ограниченной памяти & SVD раз в $N$ шагов \\
\bottomrule
\end{tabular}
\end{center}

\section{Передний край}

Направления, которые стоит держать в виду, но не брать по умолчанию.

\textbf{LoRA+} (Hayou и др., 2024). Показано, что оптимальные скорости
обучения для $A$ и $B$ различаются на порядок из-за разной размерности:
$\eta_B = \lambda \eta_A$ при $\lambda \sim 16$. Одинаковая скорость,
принятая в стандартной LoRA, замедляет сходимость.

\textbf{PiSSA} (Meng и др., 2024). Вместо случайной инициализации
$A,B$ берутся главные компоненты $W_0$: $W_0 \approx U_r S_r V_r^\top$,
$A = \sqrt{S_r}V_r^\top$, $B = U_r\sqrt{S_r}$, а остаток замораживается.
Адаптер сразу начинает с наиболее значимого подпространства
и сходится быстрее.

\textbf{VeRA} (Kopiczko и др., 2024). $A$ и $B$ общие для всех слоёв
и заморожены; обучаются только два вектора масштаба на слой. Даёт
сокращение числа параметров ещё на порядок при сопоставимом качестве.

\textbf{Spectrum} и \textbf{LISA}. Оба выбирают подмножество слоёв для
полного обучения: Spectrum — по отношению сигнал/шум в спектре весов,
LISA — случайно с перевыбором каждые $K$ шагов. Промежуточный вариант
между LoRA и полным дообучением.

\section{Практика замера}

Сравнивать методы имеет смысл по четырём числам одновременно: доля
обучаемых параметров, значение функции потерь, время шага и пиковая
память. Метод, выигравший по потерям ценой пятикратного времени,
на итерациях проиграет.

Обязателен замер \emph{до} обучения: без базовой линии непонятно, что
дал метод, а что модель умела и так. Реализовано в
\qq{vlmkit.compare.sweep} — каждый прогон начинается с чистой загрузки,
иначе адаптер предыдущего метода протекает в следующий замер.

\end{document}
